%=============================================================================
% PortSCRA -- Supply-Chain Cyber Risk Assessment DSS for Small EU Ports
% IEEE conference format (IEEEtran).
%
% Compile:  pdflatex portscra_report.tex   (run twice for cross-references)
% No external .bib is required; references use thebibliography.
%=============================================================================
\documentclass[conference]{IEEEtran}
\IEEEoverridecommandlockouts

\usepackage[T1]{fontenc}
\usepackage[utf8]{inputenc}
\usepackage{cite}
\usepackage{amsmath,amssymb}
\usepackage{graphicx}
\usepackage{booktabs}
\usepackage{textcomp}
\usepackage{url}
\usepackage{xcolor}
\usepackage[hidelinks]{hyperref}

\newcommand{\eur}{\mbox{\texteuro}}

\begin{document}

\title{PortSCRA: A Quantitative Decision Support System for\\
Supply-Chain Cyber Risk in Small European Ports}

\author{\IEEEauthorblockN{Ibrahim Dodo}
\IEEEauthorblockA{\textit{Independent research project} \\
\textit{Decision Support Systems}\\
Karlskrona, Sweden \\
idodo437@gmail.com}}

\maketitle

\begin{abstract}
Small European ports face a legal obligation to manage supply-chain cyber risk
under Directive (EU) 2022/2555 (NIS2), yet the tools available to them stop at a
qualitative heat map. Commercial third-party risk platforms are enterprise-priced,
black-box and port-blind; the accessible open tools produce a risk register rather
than a decision. This paper presents PortSCRA, a decision support system that
quantifies supplier-borne cyber risk for a small port in the units the problem is
actually posed in --- euro of revenue, hours of downtime, percent of operations
disturbed --- and then selects a mitigation portfolio under a budget constraint.
The system couples a dependency graph (vendor $\rightarrow$ asset $\rightarrow$
port service) built on the ENISA port taxonomies with a hybrid vendor-likelihood
model fusing a NIS2-mapped questionnaire, AHP-elicited criterion weights and live
vulnerability intelligence (EPSS, CISA KEV, OpenSSF Scorecard). Consequences are
quantified by FAIR-style Monte Carlo simulation calibrated from four seasonal weeks
of Danish Maritime Authority AIS data and the audited accounts of the anchor port,
Rønne (Bornholm). A submodular knapsack optimiser then answers the procurement
question directly. Validation is threefold: replay of six documented port incidents
with anchoring and held-out sets kept separate, a controlled duel against a
5$\times$5 qualitative matrix, and a persona sweep. Three results are reported that
intuition did not predict: deriving the hourly outage cost from audited accounts
reorders the risk ranking so that invoice fraud overtakes ransomware; a mitigation
portfolio reduces the mean expected loss by 77\% but the P99 tail by only 63\%, so
reporting the mean overstates the protection purchased; and return on security
investment turns negative before the control catalogue is exhausted. The held-out
incident replay failed on first execution and forced a structural change to the
consequence model, which lowered the system's own loss estimate.
\end{abstract}

\begin{IEEEkeywords}
decision support systems, supply-chain security, cyber risk quantification, FAIR,
maritime cybersecurity, NIS2, port security, AIS
\end{IEEEkeywords}

%=============================================================================
\section{Introduction}

Maritime transport carries the majority of European Union trade through more than
1{,}200 seaports \cite{enisa2019}. Those ports have digitalised rapidly, and in
doing so have acquired a dependency structure that is neither well understood nor
well governed: terminal operating systems, port community systems, cloud
productivity suites, gate and crane equipment with vendor-maintained remote access,
and a long tail of freely obtained software whose provenance nobody records.

The attacks that matter increasingly arrive \emph{through} that structure rather
than at its perimeter. SolarWinds distributed trojanized, validly signed updates to
approximately 18{,}000 organisations \cite{solarwinds}. Kaseya's remote management
tooling was used to deploy ransomware to roughly 1{,}500 downstream businesses
\cite{kaseya}. A single managed file-transfer product, MOVEit, exposed data at over
2{,}700 organisations \cite{moveit}. The xz-utils backdoor (CVE-2024-3094) placed
malicious code inside a compression library present on almost every Linux system,
and was discovered by accident \cite{xzutils}. Verizon's Data Breach Investigations
Report records third-party involvement in a substantial and growing share of
breaches \cite{dbir2025}.

Ports are not spectators to this. The 2017 NotPetya incident cost Maersk an
estimated USD 250--300 million and halted 76 terminals \cite{maersk}. The Port of
Nagoya lost roughly two and a half days of container operations to LockBit
ransomware in 2023 \cite{nagoya}. DP World Australia lost approximately three days
in the same year. Antwerp's container release systems were manipulated over a
two-year period to extract narcotics containers \cite{enisa2019}.

\subsection{The decision maker}

This work targets a specific and under-served decision maker: the person
responsible for cybersecurity at a port of 10 to 100 employees, who is typically
also responsible for IT, operations, or both, and who has no security team.

Their obligation is not optional. Ports are Annex~I ``Transport --- Water''
entities under NIS2 \cite{nis2}. Article~21(2)(d) makes supply-chain security an
explicit legal requirement; Article~20 places training duties and personal
liability on the management body; Article~23 imposes 24-hour and 72-hour incident
reporting. Critically for the smallest ports, Article~2(2)(b) brings an entity into
scope \emph{regardless of size} where it is the sole provider in a Member State of
a service essential to critical societal or economic activities.

\subsection{Problem statement}

The brief this project addresses poses two questions:

\begin{itemize}
\item \textbf{Problem 1.} Which security tools, techniques and procedures should
the organisation adopt in order to mitigate the cyber risks associated with its
supplying vendors?
\item \textbf{Problem 2.} Given the likelihood that a specific vendor is
compromised, how large is the consequence for the business --- in revenue fall,
in operational disturbance, in downtime hours?
\end{itemize}

Problem~2 is an estimation problem. Problem~1 is a constrained portfolio-selection
problem, and treating it as a ranking exercise --- which is what every accessible
tool does --- fails to answer it.

\subsection{Contributions}

\begin{enumerate}
\item A dependency-graph formulation linking suppliers to port services through the
ENISA asset taxonomy, enabling scenario frequencies to be \emph{derived} from a
port's actual supplier list rather than assumed from a catalogue.
\item An hourly outage cost derived from audited port accounts and AIS-observed
traffic, cross-validated by two independent routes, replacing the placeholder
figures typical of this literature.
\item A budget-constrained control optimiser with a multiplicative (submodular)
risk composition, an eviction-based local search, and demonstrated equivalence to
brute-force optimum on tractable instances.
\item A validation methodology that separates anchoring from held-out incidents,
and a controlled comparison against the qualitative baseline the system is intended
to replace.
\item An open, reproducible implementation: 357 automated tests, five executable
notebooks, and an interactive dashboard.
\end{enumerate}

%=============================================================================
\section{Background and Related Work}

\subsection{Regulatory and domain frameworks}

ENISA's \emph{Port Cybersecurity} report \cite{enisa2019} provides the asset and
threat taxonomies used throughout this work: ten asset categories spanning fixed
and mobile infrastructure, OT systems and end-devices, IT systems, networks, safety
and security systems, information, and people. Its companion, \emph{Cyber Risk
Management for Ports} \cite{enisa2020}, defines a four-phase approach mapped to the
ISPS Code: identify cyber-related assets and services; identify and evaluate risks;
identify security measures; assess cybersecurity maturity. PortSCRA is, in
substance, a computational operationalisation of those four phases for ports too
small to perform them manually.

The IMO addresses maritime cyber risk at the level of guidance
(MSC-FAL.1/Circ.3, MSC.428(98)), and the CIS Controls v8 \cite{cis} supply the
mitigation catalogue, with Implementation Group~1 (``essential cyber hygiene'')
being the appropriate target for an organisation of this size.

\subsection{Risk assessment methodology}

Qualitative matrices remain dominant in practice, and CSET \cite{cset} and MONARC
\cite{monarc} --- the two most accessible free tools --- both terminate in one.
Cox \cite{cox2008} shows that arithmetic on ordinal severity scales is formally
unsound, producing range compression and ranking reversals. More practically for
this problem, a matrix cannot produce the euro and hour figures the brief demands.

Factor Analysis of Information Risk (FAIR) \cite{fair} decomposes risk into loss
event frequency and loss magnitude and evaluates the result by Monte Carlo
simulation, yielding distributions rather than point estimates. Hubbard and
Seiersen \cite{hubbard} provide the calibration methodology for the expert
estimates such a model requires. Saaty's Analytic Hierarchy Process \cite{saaty}
supplies criterion weights from pairwise comparisons together with a consistency
ratio that exposes incoherent judgements rather than averaging them away.

For threat identification, MITRE ATT\&CK \cite{attack} supplies technique
identifiers; STRIDE is used only as a completeness audit, since it presumes a
data-flow diagram of a system one owns --- precisely what is unavailable for a
third party --- and the overall philosophy follows the risk-centric PASTA approach
\cite{pasta} rather than a software-centric one.

Vulnerability prioritisation draws on EPSS \cite{epss}, which estimates the
probability of exploitation within 30 days and is empirically superior to severity
scoring for triage; CISA's Known Exploited Vulnerabilities catalogue \cite{kev};
and OpenSSF Scorecard \cite{scorecard} for open-source components, which is the
only assessment available where there is no vendor to question.

\subsection{Maritime cyber risk research}

Tam and Jones' MaCRA \cite{macra} is the reference model-based framework for
maritime cyber risk, but is ship- and voyage-centric, addresses neither the supply
chain nor economic consequence, and does not produce a procurement decision.
Recent work has combined AHP, FMEA and Bayesian updating for maritime
infrastructure, and an optimisation-based approach to estimating economic losses
from port cyber-attacks \cite{portloss} validates the direction of the consequence
model presented here without addressing supplier likelihood.

Visky's research programme at Tallinn University of Technology is the closest
adjacent work on the data side. His doctoral thesis addresses cybersecurity for
maritime operational technology \cite{viskythesis}, and Benterki, Visky, Vain and
Tsiopoulos \cite{benterki2025} apply incremental inductive logic programming to
learn spoofing-attack detection rules from real AIS data recorded at Tallinn
harbour. Visky also published the AIS dataset underlying that work
\cite{viskydataset}, explicitly to address the absence of publicly available,
real-environment AIS data for validating anomaly-detection research. That programme
is the authority cited for the AIS spoofing scenario in this paper
(Section~\ref{sec:scenarios}), and its dataset was formally evaluated as a
candidate primary data source (Section~\ref{sec:datachoice}).

\subsection{Research gap}

No published work combines (i) supplier-specific compromise likelihood, (ii)
port-specific consequence quantification in operational and financial units, and
(iii) prescriptive, budget-constrained mitigation selection, in a deployable system
usable by a port without a security function. PortSCRA addresses that combination.

%=============================================================================
\section{System Architecture}

The system is organised as a pipeline in which each stage consumes the previous
stage's output, so that changing a supplier changes the recommended purchase.

\begin{enumerate}
\item \textbf{Domain model.} ENISA asset and threat taxonomies, six port services,
and a scenario catalogue.
\item \textbf{Dependency graph.} Vendors supply assets; assets support services;
services carry apportioned revenue.
\item \textbf{Likelihood engine.} Questionnaire plus live signals yields a
per-vendor annual compromise probability band.
\item \textbf{Consequence engine.} FAIR-style Monte Carlo over scenarios, with the
hourly cost derived from port economics.
\item \textbf{Decision layer.} Budget-constrained portfolio selection, efficient
frontier, tornado sensitivity.
\item \textbf{Supporting instruments.} Supplier security requirements, awareness
assessment, ENISA maturity profile.
\end{enumerate}

\subsection{Scenario catalogue}
\label{sec:scenarios}

Seven scenarios are modelled. Each is constructed from three independent frames ---
an ENISA threat category, one or more MITRE ATT\&CK techniques, and a STRIDE
completeness check --- and each is anchored to a documented incident so that its
parameters are not arbitrary. Table~\ref{tab:scenarios} summarises them.

\begin{table}[t]
\caption{Scenario catalogue with framework mappings and incident anchors}
\label{tab:scenarios}
\centering
\footnotesize
\begin{tabular}{@{}llp{2.5cm}@{}}
\toprule
\textbf{ID} & \textbf{Scenario} & \textbf{Anchor} \\
\midrule
S1 & Trojanized update to TOS/PCS & SolarWinds (2020) \\
S2 & MSP compromise $\rightarrow$ ransomware & Kaseya, Nagoya, DP World \\
S3 & SaaS account takeover $\rightarrow$ invoice fraud & BEC against maritime agents \\
S4 & OT/IoT compromise (gates, cranes, CCTV) & Antwerp (2011--13) \\
S5 & AIS/GNSS spoofing or jamming & Baltic GNSS interference \\
S6 & Compromised freeware/OSS dependency & xz-utils, CVE-2024-3094 \\
S7 & Vendor data breach & MOVEit (2023) \\
\bottomrule
\end{tabular}
\end{table}

Scenario S5 warrants specific attention because AIS is unauthenticated by design: a
position report is a claim, not a measurement. Sustained GNSS interference has been
reported across the Baltic region since 2023. The detection of such manipulation is
the subject of Visky's research programme \cite{benterki2025,viskythesis}, and the
dataset published to support it \cite{viskydataset} remains the appropriate
empirical resource should this scenario require demonstration rather than
modelling.

%=============================================================================
\section{Data and Calibration}

\subsection{Anchor port selection}

The anchor port is Rønne, on the Danish island of Bornholm. The selection criteria
were that the port be genuinely small rather than a scaled-down proxy, that its
traffic be genuinely mixed, that it lie in the Baltic region as the brief framed,
and that reliable open data exist for it.

Rønne satisfies all four. It handles over 3{,}000 ship calls, 1.4~Mt of cargo and
approximately 2~million passengers annually, across ferry routes to Ystad, Køge,
Sassnitz and Świnoujście, plus cargo, cruise and offshore-wind project traffic. Its
operating company employs 34 full-time equivalents.

That headcount is analytically important. At 34 FTE the port falls \emph{below}
NIS2's medium-enterprise threshold, and would be out of scope on size alone. As the
sole supply port for an island population it is nonetheless in scope under
Article~2(2)(b). Rønne is therefore a precise instance of the organisation this
system targets: legally obligated, operationally critical, and far too small to
operate a security function. Being an island port also sharpens the consequence
model, since an outage severs a supply line for which no road alternative exists
--- a harm ENISA identifies explicitly as more severe for islands \cite{enisa2020}.

\subsection{Dataset selection}
\label{sec:datachoice}

Two candidate AIS sources were evaluated.

The Danish Maritime Authority publishes daily AIS archives covering all Danish
waters, free of charge and without registration, from 2006 to within days of the
present \cite{dma}. Files are approximately 600~MB compressed per day, with 26
documented columns.

The alternative was the dataset published by Visky \cite{viskydataset}: six months
of AIS recorded near the Port of Tallinn, released under Apache 2.0 specifically to
provide the real-environment data that AIS anomaly-detection research had lacked.
It is a purpose-built maritime-cybersecurity research resource by an author whose
doctoral work is in exactly this domain \cite{viskythesis}, and it is a genuinely
valuable contribution.

It was nonetheless rejected as the primary source, for a reason unrelated to its
quality. The Port of Tallinn handles approximately 29.8~Mt of cargo, over
2~million TEU and 7.75--10.6~million passengers across five harbours. Anchoring on
Tallinn would have contradicted the small-port premise on which the entire system
rests: a port of that scale has a security function and does not need this DSS.
A secondary consideration was that the dataset's structure is undocumented and its
volume (3.58~GB compressed, 20.4~GB expanded for a single locality over six months)
suggests message-level raw data suited to the anomaly-detection task it was built
for rather than to the port-call statistics required here. The dataset and its
associated publications are retained as the citation basis for scenario S5.

\subsection{AIS processing}

Four full Monday-to-Sunday weeks --- one per season --- were downloaded and reduced
to port calls: approximately 18~GB ingested, 82~kB retained. A port call is defined
as a maximal run of stationary reports (speed over ground below 0.5~kn, or
navigational status indicating moored) from one vessel within the port polygon,
with consecutive reports no more than 30 minutes apart and total duration at least
one hour.

Two corrections were required, both identified by comparing output against reality
rather than by code inspection.

First, the bounding box captured the marina alongside the commercial basin. The
data itself supplied the remedy: commercial berths occupy latitudes
55.0875--55.0994 and recreational craft 55.1033--55.1056, a clean separation. Before
correction, 29 of 50 vessels observed on a July day were sailing or pleasure craft,
which carry no port dues and would have distorted any per-call economic figure.

Second, and more consequentially, the truncation tolerance was tighter than AIS's
own reporting interval. ITU-R M.1371 specifies a three-minute reporting interval for
a moored Class~A vessel, so a ship already alongside at midnight frequently produced
its first message of the day at 00:02 and was recorded as a fresh arrival. The
resulting dwell distribution was pinned at 23.9 hours at the 90th percentile --- an
artefact of the observation window rather than a berth pattern. The tolerance was
widened to five minutes and calls split by the daily file boundary were stitched.

The effect is substantial. The Passenger median dwell fell from 9.01 to 2.70 hours,
which is a realistic ro-pax turnaround and matches directly observed HAMMERSHUS
calls; 66 multi-day stays, previously invisible, were recovered.

\subsection{Reconciliation}

The AIS-derived call count was reconciled against Danmarks Statistik's official
series for Rønne (table SKIB101) \cite{dst}. On the corrected basis, commercial
calls annualise to 2{,}203 and all-vessel calls to 4{,}993. The official figures ---
3{,}043 in 2024 and 2{,}933 in 2025 --- fall \emph{inside} that bracket. Bracketing
is a stronger result than a single ratio: the official series evidently counts more
than the strict commercial set (including fishing, tugs and service craft) but
excludes the recreational traffic the all-vessel count captures.

Seasonality is pronounced: commercial calls per day range from 4.0 in winter to 7.9
in summer, a 1.96$\times$ swing, with autumn nearly matching summer owing to
offshore-wind installation campaigns. A consequence model calibrated on July alone
would overstate January risk by approximately half.

\subsection{Hourly outage cost}

The cost of one hour of disrupted operations is derived from Rønne Havn A/S's 2024
audited accounts \cite{ronne2024}: net revenue DKK~169.8~million
($\approx$~\eur{}22.76~million), EBITDA DKK~101.7~million, 3{,}043 calls, 34~FTE.
Two reductions are applied, both modelled as distributions because both are
uncertain:

\begin{itemize}
\item \textbf{Call-dependent share}, PERT(0.45, 0.60, 0.80). A cyber outage stops
vessel handling but does not stop the clock on land leases and project-area rental,
which are billed on time and grew sharply at Rønne in 2024. No segment reporting is
available --- the income statement is classified by nature --- so this is an
estimate, documented as such.
\item \textbf{Unrecoverable share}, PERT(0.20, 0.40, 0.70). Ferry and bulk traffic
largely catches up after restoration, as both Nagoya and DP World demonstrated.
What is genuinely lost is diverted traffic, cruise calls that cannot be
rescheduled, and closed offshore project windows.
\end{itemize}

The result is \eur{}659 per hour (P10 \eur{}451, P90 \eur{}883) against 8{,}760
service hours. Rønne is treated as a 24-hour operation because Bornholmslinjen
sails overnight, although 67\% of commercial arrivals fall between 06:00 and 22:00.

This figure was cross-validated by an independent derivation --- revenue per call
multiplied by the observed call rate --- which agrees to a ratio of 1.00 on a
matched call basis. The cross-check justified itself immediately: it exposed a
15$\times$ unit error in which daily revenue was converted to hourly using a factor
of 24 rather than $365/\textrm{service hours}$.

%=============================================================================
\section{Vendor Compromise Likelihood}

\subsection{Hybrid evidence model}

Neither available evidence source is sufficient alone for a small port's supplier
list, which mixes global SaaS providers with two-person integrators. A questionnaire
reaches every supplier but is self-reported and static; public signals are objective
and current but exist only where there is published software with assigned CVEs.
The two are therefore fused.

The questionnaire comprises 20 questions across seven criteria, each mapped to
points of NIS2 Article~21(2), so that administering it simultaneously produces the
supply-chain diligence evidence the directive requires. It is deliberately short:
a small port has no leverage to send a 300-question industry-standard questionnaire
to its crane OEM and expect a response.

The seven criteria are access and privilege, asset criticality, vendor security
posture, update-channel integrity, incident history and transparency, port data
held, and substitutability. Weights are obtained by AHP pairwise comparison with a
consistency check. An elicitation session produced a consistency ratio of 0.044 ---
comfortably within Saaty's 0.10 threshold --- and weights closely tracking
literature-derived defaults (access 0.238 against 0.24), which is mild
corroboration of both.

\subsection{Live vulnerability intelligence}

External signals are fetched from three free sources: EPSS \cite{epss} for
exploitation probability, CISA KEV \cite{kev} for observed exploitation, and
OpenSSF Scorecard \cite{scorecard} for open-source repository hygiene. Responses
are cached with a time-to-live, and a feed that cannot be reached leaves the signal
\emph{unset} and reports the failure, never returning a value that resembles a
clean result.

Signal weighting is deliberately asymmetric: presence in KEV raises assessed risk
substantially, while absence lowers it only slightly. The live data supplies a
compelling justification for this. CVE-2024-3094 --- the xz-utils backdoor --- scores
0.86 on EPSS yet is \emph{absent} from KEV, because it was discovered before it was
exploited in the wild. One of the most serious supply-chain compromises yet
documented would score clean on a KEV-only check.

\subsection{Scoring}

Criterion scores are combined using the AHP weights, renormalised over the criteria
actually answered, producing a posture score in $[0,1]$. Posture maps to an exposure
multiplier by geometric interpolation over $[0.4, 2.5]$, since risk multipliers
compound; the range is bounded because no questionnaire justifies asserting that a
supplier is an order of magnitude safer than average. The product of exposure and
signal multipliers scales a documented base rate, PERT(0.02, 0.05, 0.12), into a
per-vendor probability band.

Two design decisions merit statement. A vendor that supplies no evidence scores
0.5, not 0: if silence produced a favourable score, declining to answer would be
the cheapest route to appearing safe. And partial responses are scored only on what
was answered, so a supplier answering the access questions badly is not credited
for the questions it skipped.

%=============================================================================
\section{Risk Quantification}

\subsection{The dependency graph}

The join that makes the system port-specific rather than generic maps vendors to
the ENISA asset classes they supply, and asset classes to the six port services
they support: vessel handling, cargo handling, passenger handling, gate operations,
declarations, and billing. Each service carries an apportioned share of
call-dependent revenue.

The graph answers a question no per-vendor score can: how much of the port does
each supplier touch? Ranked by revenue exposed rather than by security posture, the
crane OEM reaches 80\% and the managed service provider 75\%. Concentration is a
property of the topology, and a supplier touching everything is a risk even when it
is well run.

\subsection{Derived frequencies}

Scenario frequencies are constructed from the suppliers that provide the assets a
scenario targets, each contributing in proportion to its coverage of those assets.
Contributions \emph{add}, because loss event frequency is a Poisson rate; treating
them as probabilities would be incorrect. A scenario with no matching supplier
retains its generic prior rather than falling to zero, since otherwise omitting a
vendor from the register would be the cheapest way to reduce measured risk.

The derived frequencies differ materially from the priors, and in both directions:
the trojanized-update scenario is 2.2$\times$ more likely than the generic prior for
this port, the OT compromise 1.9$\times$, while SaaS takeover falls by a factor of
three and vendor data breach by four.

\subsection{Monte Carlo}

Annual loss is simulated as
\begin{equation}
L_{\textrm{annual}} = \sum_{i=1}^{N} M_i, \qquad N \sim \textrm{Poisson}(\lambda),
\end{equation}
where $\lambda$ is drawn per simulated year from the loss event frequency
distribution and each magnitude $M_i$ is
\begin{equation}
M_i = D_i \cdot c + R_i + S_i,
\end{equation}
with $D_i$ operational downtime hours, $c$ the hourly cost, $R_i$ response cost,
and $S_i$ a secondary loss realised with a stated probability.

The Poisson step is not incidental. Multiplying frequency by magnitude yields the
mean and destroys the shape, and the shape is the object of interest: a port with a
frequency of 0.08 per year has no average year, only quiet years and one very bad
one. In the reported configuration 45\% of simulated years incur no loss at all
while the 99th percentile exceeds \eur{}566{,}000. Downtime is additionally gated by
a disruption probability, discussed in Section~\ref{sec:validation}.

%=============================================================================
\section{Decision Layer}

\subsection{Formulation}

Problem~1 is formulated as constrained portfolio selection. Each control $j$ has an
annual cost $k_j$ and, per scenario $s$, a frequency reduction $f_{js}$ and a
magnitude reduction $m_{js}$ --- kept separate because backups do not prevent
ransomware, they make it survivable, and a single ``effectiveness'' scalar cannot
express that difference. Residual expected loss is
\begin{equation}
\textrm{ALE}(X) = \sum_{s} \textrm{ALE}_s
\prod_{j \in X}(1 - f_{js}) \prod_{j \in X}(1 - m_{js}),
\end{equation}
minimised subject to $\sum_{j \in X} k_j \le B$.

Composition is multiplicative rather than additive. Two controls each halving a
frequency leave a quarter of it, not zero; additive composition would permit the
model to ``eliminate'' a risk by stacking mediocre controls and would systematically
over-purchase. Multiplicativity also renders the objective submodular, which is what
makes greedy selection principled rather than merely convenient.

\subsection{Search}

Selection proceeds greedily by marginal reduction per euro, with marginals
recomputed after each choice so that overlap is accounted for. Greedy alone,
however, falls into the classical knapsack trap: at a \eur{}1{,}000 budget it
purchased two inexpensive high-ratio controls totalling \eur{}800 when a single
\eur{}1{,}000 control was strictly superior, and no one-for-one swap could recover
the position because affording the better control required evicting both. An
eviction-based local search was added, admitting a candidate by removing incumbents
in ascending order of marginal contribution. The augmented search matches
brute-force optimum at every budget tested.

The control catalogue comprises 20 implementable projects rather than individual
safeguards, each mapped to CIS v8 safeguards and NIS2 points. CIS v8 defines 153
safeguards and Implementation Group~1 alone contains 56; a 34-person port will not
implement 56 discrete items, and a catalogue at that granularity produces an answer
nobody acts upon.

%=============================================================================
\section{Supporting Instruments}

Three further components complete ENISA's four-phase cycle and the brief's stated
DSS tasks.

\textbf{Supplier security requirements.} Fourteen requirements are scoped by the
asset classes a supplier provides, tiered essential/recommended/advanced, and each
carries contract text, an evidence test, and the NIS2 point it discharges.
Requirements are keyed to the questionnaire questions that verify them, so gap
analysis falls out of answers already collected. A requirement with no verifying
answer is reported as \emph{unverified} rather than met: not asking is not passing.
The essential tier is capped, because a list a port is willing to lose a supplier
over must be short enough to survive a procurement negotiation.

\textbf{Awareness assessment.} The brief asks to identify awareness of
\emph{specific} vendor-related risks, so the instrument is not a fixed quiz:
vignettes are selected in the order the risk engine ranked the scenarios. The
measured score feeds back into scenario frequencies, weighted by how much of each
scenario is plausibly under human control --- invoice fraud at 0.60, a signed
trojanized update at 0.05. A tool claiming that training reduces every risk equally
would be selling training rather than measuring it. Management scores are reported
separately, because Article~20(2) obliges the management body itself to train, and
an organisational average can conceal a board that has not.

\textbf{Maturity.} ENISA phase four is implemented using the three levels defined in
\cite{enisa2020}: Basic, Intermediate, Optimal. An optimised portfolio implies a
maturity profile, allowing a budget to be read in ENISA's terms as well as in euro.
The mapping deliberately caps what purchasing alone achieves: the full catalogue
reaches 2.42 of 3, because level~3 requires oversight, testing and continuous
improvement, and ENISA is explicit that maturity cannot be attained through a
checklist mentality.

%=============================================================================
\section{Validation}
\label{sec:validation}

\subsection{Incident replay}

Six documented port incidents were replayed against the model, separated into those
that \emph{anchored} the estimates and those \emph{held out}. Only the held-out set
constitutes evidence; checking a model against the data used to construct it is a
consistency check, and reporting it as validation would be dishonest.

The first execution failed two of four held-out cases, and the failure was the most
productive result of the project. The Port of Barcelona and the Port of San Diego
(both 2018) suffered ransomware without halting quay operations: IT systems were
affected while vessel operations continued. The model's downtime estimate had a
floor of 24 hours and therefore could not produce that outcome at any percentile.
It was asserting that supplier-borne ransomware invariably becomes an operational
outage, which the record contradicts.

The remedy was structural rather than a parameter adjustment. A
\emph{disruption probability} was introduced, separating ``an incident occurred''
from ``operations stopped''; simply lowering the floor would have conflated two
distinct phenomena. It was set to 0.67 for S2 from the evidence --- four of six
documented incidents disrupted operations --- and 0.50 for S1, since few of the
approximately 18{,}000 organisations that installed the SolarWinds backdoor were
actively exploited, and the resulting harm was predominantly espionage rather than
outage.

All four held-out incidents now fall inside the predicted band, and the change
\emph{lowered} the model's own expected-loss estimate for S2 from \eur{}27{,}852 to
\eur{}25{,}714. Because band containment becomes a weak test once the band begins at
zero, a sharper check is also reported: the model predicts that 33\% of ransomware
events cause no operational downtime, against 33\% (two of six) observed.

\subsection{Baseline comparison}

The system was compared against a 5$\times$5 qualitative matrix, both methods fed
the \emph{same} underlying estimates so that the comparison isolates method rather
than inputs. The result is more informative than a simple defeat of the baseline.

Rank agreement between the matrix score and the quantitative expected loss is
$+0.91$ (Spearman): the matrix \emph{orders} risks perfectly sensibly. Were that
not so, outperforming it would prove little. Where it fails is allocation. Six of
seven scenarios fall into a single ``Medium'' band spanning a factor of 4.2 in
expected loss --- Cox's range compression \cite{cox2008} made concrete --- so a heat
map cannot prioritise among them.

Spending an identical budget by each method, the optimiser wins at every level
tested, by \eur{}1{,}644 to \eur{}18{,}856 per year, peaking at $+17\%$ of baseline
at a \eur{}15{,}000 budget (Table~\ref{tab:duel}). The precise claim is therefore
that \emph{a matrix can rank well and still allocate badly}, because it cannot
perceive that an inexpensive control touching four amber risks outperforms an
expensive one touching a single red.

\begin{table}[t]
\caption{Residual expected loss, same budget spent two ways}
\label{tab:duel}
\centering
\footnotesize
\begin{tabular}{@{}rrrr@{}}
\toprule
\textbf{Budget} & \textbf{Heat-map led} & \textbf{Optimised} & \textbf{Advantage} \\
\midrule
\eur{}10,000 & \eur{}58,337 & \eur{}52,335 & \eur{}6,002 \\
\eur{}15,000 & \eur{}57,709 & \eur{}38,853 & \eur{}18,856 \\
\eur{}25,000 & \eur{}32,344 & \eur{}25,384 & \eur{}6,960 \\
\eur{}40,000 & \eur{}22,793 & \eur{}18,561 & \eur{}4,232 \\
\eur{}60,000 & \eur{}15,804 & \eur{}14,160 & \eur{}1,644 \\
\bottomrule
\end{tabular}
\end{table}

\subsection{Persona sweep}

The brief specifies 10, 50 and 100-employee personas, which bracket Rønne's actual
34~FTE. Both revenue and response costs scale with headcount, giving hourly outage
costs of \eur{}194, \eur{}969 and \eur{}1{,}937 respectively.

At a fixed \eur{}25{,}000 budget all three personas select an identical portfolio.
This is an artefact rather than a finding: \eur{}25{,}000 represents 88\% of the
smallest persona's entire annual expected loss and 19\% of the largest's. At
size-proportional budgets the three select three distinct portfolios, and the
smallest is the most instructive --- a ten-person port with \eur{}2{,}500 is directed
to purchase exactly one control: out-of-band verification of payment changes, a
\eur{}1{,}500 procedural measure that reduces the highest-ranked scenario by three
quarters. That is a credible answer for an organisation that cannot afford a
security programme at all.

%=============================================================================
\section{Results and Discussion}

\subsection{Portfolio selection}

At a \eur{}25{,}000 budget the optimiser reduces expected annual loss from
\eur{}111{,}186 to \eur{}25{,}384, a 77\% reduction at a return on security
investment of $+2.4\times$. Seven of the eight selected controls belong to CIS
Implementation Group~1. This is a \emph{result} rather than a constraint: nothing
restricted the search to IG1, which simply purchased inexpensive fundamentals
first. It constitutes independent support for CIS's own position on what carries
most value for small organisations.

\subsection{Three findings that intuition did not predict}

\textbf{The derived cost reorders the risks.} Under a placeholder hourly cost, MSP
ransomware ranked first. Once the figure was derived from audited accounts ---
\eur{}659, because a small port's revenue per hour is genuinely modest ---
SaaS-mediated invoice fraud overtook it. The reason is structural: downtime-driven
losses scale with a modest hourly revenue, whereas a fraudulent payment is direct,
immediate and unrecoverable. A qualitative assessment would in all likelihood have
rated ransomware ``critical'' and business email compromise ``medium''.

\textbf{The tail is stickier than the mean.} The selected portfolio reduces mean
expected loss by 77\% but the 99th percentile by only 63\%. The portfolio
predominantly purchases \emph{frequency} reduction --- loss-free years rise from 34\%
to 52\% --- while the bad year that still occurs remains expensive, because the
residual tail is dominated by whichever scenarios the budget could not address.
Reporting only the mean therefore \emph{overstates} the protection purchased. This
runs contrary to the usual argument for quantification, which presumes benefit
concentrates in the tail.

\textbf{Return on investment turns negative before the catalogue is exhausted.} The
efficient frontier yields ROSI of $+5.5\times$ at \eur{}9{,}000, $+2.0\times$ at
\eur{}30{,}000, $+1.0\times$ at \eur{}48{,}000, and $-0.3\times$ at full-catalogue
spend. Purchasing everything is worse than purchasing the correct subset --- a
conclusion a compliance checklist structurally cannot reach, and the strongest
single argument in this work for optimisation over enumeration.

\subsection{Methodological observation}

Two conclusions in this project reversed under a parameter sweep after appearing
robust at a single value. The first was the tail-versus-mean relationship; the
second was a comparison of generic against supplier-derived frequencies, which
appeared to produce identical portfolios at the default budget and in fact differs
at five of eight budget levels. A comparison performed at one parameter value is an
anecdote. This bears directly on the single-case-study methodology common in
applied DSS reporting.

%=============================================================================
\section{Limitations}

\textbf{Estimate provenance.} The weakest inputs are control effectiveness and the
two economic share parameters. Every estimate carries a documented source string,
and an automated test fails if any lacks one, but they remain estimates. The
tornado analysis is the appropriate response: it identifies which estimates would
change the decision if wrong, and which would not.

\textbf{Linear persona scaling.} Revenue is scaled linearly with headcount, a crude
proxy for a capital-intensive industry in which revenue per employee varies widely
with automation and business mix.

\textbf{Coverage weighting.} A vendor's contribution to a scenario is proportional
to the fraction of that scenario's target assets it supplies. This is a proxy, not a
measurement, and is the assumption in the dependency graph most deserving of
challenge.

\textbf{Single anchor port.} All calibration derives from one port. The method
transfers; the numbers do not.

\textbf{Illustrative supplier list.} The worked example uses illustrative supplier
\emph{roles}, not real companies. Naming real firms would assert claims about their
security posture for which no evidence exists. Correspondingly, no signal value is
entered by hand: an earlier iteration typed EPSS scores and KEV counts directly, and
this inflated one supplier's assessed risk from 7.6\% to 13.3\% on invented
evidence. A test now fails if any hand-entered feed value reappears.

%=============================================================================
\section{Conclusion and Future Work}

PortSCRA demonstrates that supplier-borne cyber risk for a small port can be
quantified in the units the problem is posed in, and that the resulting figures
support a materially better procurement decision than the qualitative baseline they
would replace --- by \eur{}1{,}644 to \eur{}18{,}856 per year at identical budgets.
All four ENISA risk-management phases are implemented, and every deliverable the
originating brief specifies is realised in code.

The results that matter most are those that contradicted expectation. Deriving the
hourly outage cost from audited accounts reordered the risk ranking. A mitigation
portfolio moves the mean considerably further than the tail, so the mean overstates
the protection obtained. Return on investment turns negative before the catalogue is
exhausted. And a held-out validation set failed on first execution, forcing a
structural correction that lowered the system's own loss estimate --- which is what
a validation suite is for. A validation suite that never fails is not one.

Future work falls into three directions. First, discrete-event simulation of outage
propagation, capturing the non-linearity by which a 48-hour outage costs more than
48 hours of revenue through queueing and recovery effects. Second, extension to
multiple anchor ports to test whether the calibration method transfers as the
architecture is intended to. Third, empirical treatment of the AIS spoofing
scenario using the real-environment dataset published for that purpose
\cite{viskydataset}, which would replace a modelled likelihood with a measured one.

%=============================================================================
\section*{Reproducibility}

The implementation comprises approximately 12{,}900 lines of Python across 28
modules, with 357 automated tests that require no network access, five executable
notebooks with committed outputs, and an interactive dashboard. Vulnerability
intelligence is fetched live and cached; AIS processing is reproducible from the
public archive; all derived port-call data required to reproduce the analysis is
included at 82~kB.

\section*{Statement on the Use of AI}

This work was produced with AI collaboration (Level~3). An AI assistant was used
for implementation, literature search, code review and drafting. All architectural
decisions, the methodology, the validation design and the interpretation of results
are the author's own, and every quantitative claim in this paper was verified
against the executing implementation. Several of the findings reported here ---
notably the incident-replay failure and the parameter-sweep reversals --- were
identified by tests and analyses the author specified for the purpose of finding
such errors.

%=============================================================================
\begin{thebibliography}{00}

\bibitem{enisa2019}
ENISA, ``Port Cybersecurity: Good practices for cybersecurity in the maritime
sector,'' European Union Agency for Cybersecurity, Nov. 2019.

\bibitem{enisa2020}
ENISA, ``Cyber Risk Management for Ports: Guidelines for cybersecurity in the
maritime sector,'' European Union Agency for Cybersecurity, Dec. 2020.
ISBN 978-92-9204-403-9, doi: 10.2824/671060.

\bibitem{nis2}
European Parliament and Council, ``Directive (EU) 2022/2555 concerning measures for
a high common level of cybersecurity across the Union (NIS2),'' \emph{Official
Journal of the European Union}, L333, pp. 80--152, Dec. 2022.

\bibitem{cox2008}
L. A. Cox Jr., ``What's Wrong with Risk Matrices?'' \emph{Risk Analysis}, vol. 28,
no. 2, pp. 497--512, 2008, doi: 10.1111/j.1539-6924.2008.01030.x.

\bibitem{fair}
The Open Group, ``Risk Analysis (O-RA), Version 2.0.1,'' Open Group Standard, 2021.

\bibitem{hubbard}
D. W. Hubbard and R. Seiersen, \emph{How to Measure Anything in Cybersecurity
Risk}, 2nd ed. Hoboken, NJ: Wiley, 2023.

\bibitem{saaty}
T. L. Saaty, \emph{The Analytic Hierarchy Process: Planning, Priority Setting,
Resource Allocation}. New York: McGraw-Hill, 1980.

\bibitem{attack}
MITRE, ``MITRE ATT\&CK: Enterprise and ICS Matrices.'' [Online]. Available:
\url{https://attack.mitre.org/}

\bibitem{pasta}
T. UcedaVélez and M. M. Morana, \emph{Risk Centric Threat Modeling: Process for
Attack Simulation and Threat Analysis}. Hoboken, NJ: Wiley, 2015.

\bibitem{epss}
J. Jacobs, S. Romanosky, B. Edwards, I. Adjerid, and M. Roytman, ``Exploit
Prediction Scoring System (EPSS),'' \emph{Digital Threats: Research and Practice},
vol. 2, no. 3, pp. 1--17, 2021, doi: 10.1145/3436242. [Online]. Available:
\url{https://www.first.org/epss/}

\bibitem{kev}
CISA, ``Known Exploited Vulnerabilities Catalog.'' [Online]. Available:
\url{https://www.cisa.gov/known-exploited-vulnerabilities-catalog}

\bibitem{scorecard}
Open Source Security Foundation, ``OpenSSF Scorecard.'' [Online]. Available:
\url{https://securityscorecards.dev/}

\bibitem{cis}
Center for Internet Security, ``CIS Critical Security Controls, Version 8,'' 2021.

\bibitem{macra}
K. Tam and K. Jones, ``MaCRA: a model-based framework for maritime cyber-risk
assessment,'' \emph{WMU Journal of Maritime Affairs}, vol. 18, pp. 129--163, 2019,
doi: 10.1007/s13437-019-00162-2.

\bibitem{portloss}
Y. Ben Farah, et al., ``Estimating economic losses from cyber-attacks on shipping
ports: An optimization-based approach,'' \emph{Transportation Research Part C:
Emerging Technologies}, vol. 133, 103423, 2021,
doi: 10.1016/j.trc.2021.103423.

\bibitem{viskythesis}
G. Visky, ``Cybersecurity for Maritime Operational Technology: Challenges,
Considerations and Solutions,'' Ph.D. dissertation, Dept. Software Science, Tallinn
Univ. of Technology, Tallinn, Estonia, 2025.

\bibitem{benterki2025}
A. S. Benterki, G. Visky, J. Vain, and L. Tsiopoulos, ``Using Incremental Inductive
Logic Programming for Learning Spoofing Attacks on Maritime Automatic
Identification System Data,'' in \emph{Maritime Cybersecurity}, S. Bauk, Ed. Cham:
Springer, 2025, ch. 7, doi: 10.1007/978-3-031-87290-7\_7.

\bibitem{viskydataset}
G. Visky, ``AIS Dataset: 6-month long AIS data recorded near by port of Tallinn,''
Kaggle, 2024, Apache License 2.0. [Online]. Available:
\url{https://www.kaggle.com/datasets/gaborvisky/ais-dataset}

\bibitem{dma}
Danish Maritime Authority, ``AIS Data.'' [Online]. Available:
\url{http://aisdata.ais.dk}

\bibitem{dst}
Statistics Denmark, ``SKIB101: Call of vessels, passengers and throughput of goods
in traffic ports by seaport, unit and time.'' [Online]. Available:
\url{https://www.statbank.dk/SKIB101}

\bibitem{ronne2024}
Rønne Havn A/S, ``Årsrapport 2024,'' CVR no. 27932150, Rønne, Denmark, 2025.

\bibitem{dbir2025}
Verizon, ``2025 Data Breach Investigations Report,'' Verizon Business, 2025.

\bibitem{maersk}
A.P. Møller--Maersk, ``Interim Report Q2 2017,'' Copenhagen, Denmark, Aug. 2017.

\bibitem{nagoya}
Dragos, ``OT Cybersecurity Breach Disrupts Operations at the Port of Nagoya,
Japan,'' Jul. 2023. [Online]. Available:
\url{https://www.dragos.com/blog/}

\bibitem{solarwinds}
CISA, ``Emergency Directive 21-01: Mitigate SolarWinds Orion Code Compromise,''
Dec. 2020.

\bibitem{kaseya}
CISA, ``Kaseya VSA Supply-Chain Ransomware Attack,'' Alert AA21-183A, Jul. 2021.

\bibitem{moveit}
CISA, ``\#StopRansomware: CL0P Ransomware Gang Exploits MOVEit Vulnerability,''
Alert AA23-158A, Jun. 2023.

\bibitem{xzutils}
National Vulnerability Database, ``CVE-2024-3094: Malicious code in xz-utils,''
Mar. 2024.

\bibitem{cset}
CISA, ``Cyber Security Evaluation Tool (CSET).'' [Online]. Available:
\url{https://github.com/cisagov/cset}

\bibitem{monarc}
Luxembourg House of Cybersecurity, ``MONARC: Method for an Optimised Analysis of
Risks by CASES.'' [Online]. Available: \url{https://www.monarc.lu/}

\end{thebibliography}

\end{document}
